\documentclass{beamer}

\usepackage[T2A]{fontenc}
\usepackage[utf8]{inputenc}
\usepackage[english,russian]{babel}
\input{settings.tex}


\title{Метод внутренней точки}
\subtitle{Вычислительный интеллект, Лекция  14}
\author{Сергей Савин}
\centering
\date{\mydate}



\begin{document}
\maketitle


%\begin{frame}{Content}
%
%\begin{itemize}
%	\item Newton method
%	\item Barrier functions
%	\item Interior point method
%	\item Analytic center of linear inequalities
%\end{itemize}
%
%\end{frame}

\begin{frame}{Квадратичная программа с одним неравенством}
	\begin{flushleft}
		
		Рассматриваем:
		\begin{align*}
			\min_{\bo{x}} \quad & \tfrac{1}{2} \bo{x}^\top \bo{Q} \bo{x} + \bo{q}^\top \bo{x} \\
			\text{s.t.} \quad & \bo{A}\bo{x} = \bo{b} \\
			& \bo{c}^\top \bo{x} \le d
		\end{align*}
		
		\vspace{0.5em}
		
		Вводим:
		\begin{itemize}
			\item Переменная ослабления: $z \ge 0$, такая что $\bo{c}^\top \bo{x} + z = d$
			\item Двойственные переменные:
			\begin{itemize}
				\item $\bo{y}$ для ограничения-равенства
				\item $s \ge 0$ для ограничения-неравенства
			\end{itemize}
		\end{itemize}
		
		\vspace{0.5em}
		
		Переменные решения: $(\bo{x}, \bo{y}, z, s)$
		
	\end{flushleft}
\end{frame}

\begin{frame}{Условия ККТ (Релаксированные)}
	\begin{flushleft}
		
		Условия Каруша-Куна-Таккера (KKT):
		
		\vspace{0.5em}
		
		\textbf{Стационарность:}
		\[
		\bo{Q}\bo{x} + \bo{q} + \bo{A}^\top \bo{y} + s \bo{c} = 0
		\]
		
		\textbf{Прямая допустимость:}
		\[
		\bo{A}\bo{x} = \bo{b}, \quad \bo{c}^\top \bo{x} + z = d
		\]
		
		\textbf{Релаксированная дополняющая нежесткость:}
		\[
		z s = \mu
		\]
		
		\textbf{Положительность:}
		\[
		z > 0, \quad s > 0
		\]
		
		\vspace{0.5em}
		
		$\mu > 0$ — параметр барьера.
		
	\end{flushleft}
\end{frame}

\begin{frame}{Нелинейная система}
	\begin{flushleft}
		
		Определяем невязки:
		\begin{align*}
			\bo{r}_d &= \bo{Q}\bo{x} + \bo{q} + \bo{A}^\top \bo{y} + s \bo{c} \\
			\bo{r}_p^{(1)} &= \bo{A}\bo{x} - \bo{b} \\
			\bo{r}_p^{(2)} &= \bo{c}^\top \bo{x} + z - d \\
			r_c &= z s - \mu
		\end{align*}
		
		\vspace{0.5em}
		
		Решаем нелинейную систему:
		\[
		\bo{r}_d = 0, \quad \bo{r}_p^{(1)} = 0, \quad \bo{r}_p^{(2)} = 0, \quad r_c = 0
		\]
		
	\end{flushleft}
\end{frame}

\begin{frame}{Линеаризация Ньютона}
	\begin{flushleft}
		
		Пусть $(\Delta \bo{x}, \Delta \bo{y}, \Delta z, \Delta s)$ — шаг Ньютона.
		
		\vspace{0.5em}
		
		Линеаризуя систему:
		
		\vspace{0.5em}
		
		\begin{align*}
			\bo{Q} \Delta \bo{x} + \bo{A}^\top \Delta \bo{y} + \bo{c} \, \Delta s &= -\bo{r}_d \\
			\bo{A} \Delta \bo{x} &= -\bo{r}_p^{(1)} \\
			\bo{c}^\top \Delta \bo{x} + \Delta z &= -\bo{r}_p^{(2)} \\
			s \, \Delta z + z \, \Delta s &= -r_c
		\end{align*}
		
	\end{flushleft}
\end{frame}

\begin{frame}{Выбор длины шага}
	\begin{flushleft}
		
		Обновление:
		\[
		z^+ = z + \alpha \Delta z, \quad s^+ = s + \alpha \Delta s
		\]
		
		\vspace{0.5em}
		
		Выбираем длину шага для сохранения положительности:
		\[
		z^+ > 0, \quad s^+ > 0
		\]
		
	\end{flushleft}
\end{frame}

\begin{frame}{Обновление барьера и остановка}
	\begin{flushleft}
		
		Обновление параметра барьера:
		\[
		\mu = \sigma (z s), \quad \sigma \in (0,1)
		\]
		
		\vspace{0.5em}
		
		Интерпретация:
		\begin{itemize}
			\item $zs$ измеряет зазор дополняющей нежесткости
			\item $\mu \to 0$ обеспечивает оптимальность
		\end{itemize}
		
		\vspace{0.5em}
		
		Условия остановки:
		\begin{itemize}
			\item Прямая допустимость: $\|\bo{A}\bo{x} - \bo{b}\|$ мала
			\item Допустимость неравенства: $|\bo{c}^\top \bo{x} + z - d|$ мала
			\item Двойственная допустимость: $\|\bo{Q}\bo{x} + \bo{q} + \bo{A}^\top \bo{y} + s\bo{c}\|$ мала
			\item Дополняющая нежесткость: $zs$ мала
		\end{itemize}
		
	\end{flushleft}
\end{frame}









\myqrframe



\end{document}
